%% ─────────────────────────────────────────────────────────────────────────────
\section{The Complete API Specification}
%% ─────────────────────────────────────────────────────────────────────────────

This section gives you complete function contracts. A \emph{contract} specifies: what you must provide (preconditions), what the function guarantees when it returns (postconditions), and what it guarantees on both success and failure paths.

\subsection{The data structure (\texttt{include/instance.h})}

\begin{lstlisting}[style=cstyle]
#ifndef TSP_INSTANCE_H
#define TSP_INSTANCE_H

/* ── TSPInstance ────────────────────────────────────────────────────────────
 *
 *  Holds a single TSP instance loaded from a TSPLIB coordinate file.
 *
 *  Fields:
 *    n         Number of cities. Always >= 1 when the struct is valid.
 *              Zero indicates an uninitialized or freed struct.
 *
 *    coords_x  Pointer to an array of n x-coordinates, allocated on the heap.
 *              coords_x[i] is the x-coordinate of city i (0-indexed).
 *              NULL if not yet allocated.
 *
 *    coords_y  Pointer to an array of n y-coordinates, allocated on the heap.
 *              coords_y[i] is the y-coordinate of city i (0-indexed).
 *              NULL if not yet allocated.
 *
 *    dist      Pointer to a flat row-major n*n distance matrix, heap-allocated.
 *              dist[i*n + j] is the Euclidean distance from city i to city j.
 *              NULL if not yet built.
 *
 *  Invariants (when valid):
 *    dist[i*n + i] == 0.0          for all i
 *    dist[i*n + j] == dist[j*n+i]  for all i, j  (symmetry)
 *    dist[i*n + j] >= 0.0          for all i, j
 *
 *  Ownership:
 *    The struct owns all three pointers. tsp_instance_free() must be called
 *    exactly once on any struct touched by tsp_instance_load(), whether or
 *    not the load succeeded.
 * ─────────────────────────────────────────────────────────────────────────── */
typedef struct {
    int     n;
    double *coords_x;
    double *coords_y;
    double *dist;
} TSPInstance;

/* Access macro. i and j must be in [0, inst->n). */
#define DIST(inst, i, j)  ((inst)->dist[(i) * (inst)->n + (j)])

/* ── Function declarations ───────────────────────────────────────────────── */

int  tsp_instance_load(const char *path, TSPInstance *out);
void tsp_instance_free(TSPInstance *inst);
int  tsp_build_distance_matrix(TSPInstance *inst);

#endif /* TSP_INSTANCE_H */
\end{lstlisting}

\newpage

\begin{apispec}{tsp\_instance\_load}

\textbf{Signature:}
\begin{lstlisting}[style=cstyle, numbers=none]
int tsp_instance_load(const char *path, TSPInstance *out);
\end{lstlisting}

\textbf{Arguments:}
\begin{itemize}
  \item \texttt{path} --- A null-terminated C string giving the file system path to a TSPLIB-format \texttt{.tsp} file. Must not be \texttt{NULL}.
  \item \texttt{out} --- Pointer to a \texttt{TSPInstance} that the caller provides. The function writes the loaded instance into this struct. Must not be \texttt{NULL}. The struct's contents before the call are ignored; the function overwrites them.
\end{itemize}

\textbf{Returns:}
\begin{itemize}
  \item \texttt{0} on success. \texttt{out} is fully populated: \texttt{n} $\geq$ 1, \texttt{coords\_x} and \texttt{coords\_y} are heap-allocated arrays of length \texttt{n}, \texttt{dist} is \texttt{NULL} (matrix is not built by this function).
  \item \texttt{-1} on any failure. Failure conditions include: file not found, file is empty, \texttt{NODE\_COORD\_SECTION} is missing, any coordinate line is malformed, any allocation fails. On failure, \texttt{out} is left in a state that is safe to pass to \texttt{tsp\_instance\_free()}: all pointers it owns are either valid or \texttt{NULL}.
\end{itemize}

\textbf{What it does (step by step):}
\begin{enumerate}[label=\arabic*.]
  \item Initialise \texttt{*out} to the zero struct (\texttt{n=0}, all pointers \texttt{NULL}). This ensures \texttt{free} is safe even if we return early.
  \item Open the file at \texttt{path}. Return \texttt{-1} if it cannot be opened.
  \item Scan lines until \texttt{DIMENSION: N} (or \texttt{DIMENSION : N}) is found. Parse $N$ as the city count. Return \texttt{-1} if not found before \texttt{NODE\_COORD\_SECTION} or EOF.
  \item Allocate \texttt{coords\_x[N]} and \texttt{coords\_y[N]}. Return \texttt{-1} if either allocation fails.
  \item Scan lines until \texttt{NODE\_COORD\_SECTION} is found. Return \texttt{-1} if not found.
  \item For each of the next $N$ lines, parse with \texttt{sscanf(line, "\%d \%lf \%lf", \&id, \&x, \&y)}. If any line does not produce exactly 3 fields, return \texttt{-1}.
  \item Store \texttt{x} in \texttt{coords\_x[id-1]} and \texttt{y} in \texttt{coords\_y[id-1]}. (TSPLIB IDs are 1-indexed; convert to 0-indexed.)
  \item Close the file. Set \texttt{out->n = N}. Return \texttt{0}.
\end{enumerate}

\textbf{What it does NOT do:}
\begin{itemize}
  \item It does \emph{not} build the distance matrix. Call \texttt{tsp\_build\_distance\_matrix} separately.
  \item It does \emph{not} validate that city IDs form a contiguous sequence starting at 1.
\end{itemize}

\end{apispec}

\vspace{8pt}

\begin{apispec}{tsp\_instance\_free}

\textbf{Signature:}
\begin{lstlisting}[style=cstyle, numbers=none]
void tsp_instance_free(TSPInstance *inst);
\end{lstlisting}

\textbf{Arguments:}
\begin{itemize}
  \item \texttt{inst} --- Pointer to a \texttt{TSPInstance}. May be \texttt{NULL} (the function is a no-op in that case). The struct it points to may be partially initialised (some pointers non-\texttt{NULL}, others \texttt{NULL}).
\end{itemize}

\textbf{Returns:} Nothing.

\textbf{What it does:}
\begin{enumerate}[label=\arabic*.]
  \item If \texttt{inst} is \texttt{NULL}, return immediately.
  \item Call \texttt{free(inst->coords\_x)} if it is not \texttt{NULL}. Set to \texttt{NULL}.
  \item Call \texttt{free(inst->coords\_y)} if it is not \texttt{NULL}. Set to \texttt{NULL}.
  \item Call \texttt{free(inst->dist)} if it is not \texttt{NULL}. Set to \texttt{NULL}.
  \item Set \texttt{inst->n = 0}.
\end{enumerate}

\textbf{Why nullify after free?} Calling \texttt{free()} on a pointer does not set the pointer to \texttt{NULL}. If anything later accidentally calls \texttt{free()} again on the same pointer, it is undefined behaviour (double-free). Setting to \texttt{NULL} after freeing makes the struct safe to call \texttt{free} on a second time (\texttt{free(NULL)} is always a no-op in C).

\textbf{Why safe on partial init?} During \texttt{tsp\_instance\_load}, several allocations happen. If the third allocation fails and you return \texttt{-1}, the struct has some non-\texttt{NULL} pointers and some \texttt{NULL} ones. If \texttt{free} only frees non-\texttt{NULL} pointers, the caller can always call \texttt{tsp\_instance\_free} unconditionally without leaking.

\end{apispec}

\vspace{8pt}

\begin{apispec}{tsp\_build\_distance\_matrix}

\textbf{Signature:}
\begin{lstlisting}[style=cstyle, numbers=none]
int tsp_build_distance_matrix(TSPInstance *inst);
\end{lstlisting}

\textbf{Arguments:}
\begin{itemize}
  \item \texttt{inst} --- Pointer to a \texttt{TSPInstance} that has already been successfully loaded: \texttt{inst->n} $\geq 1$, \texttt{coords\_x} and \texttt{coords\_y} are valid, \texttt{dist} is \texttt{NULL} (or the function will overwrite it).
\end{itemize}

\textbf{Returns:}
\begin{itemize}
  \item \texttt{0} on success. \texttt{inst->dist} points to a heap-allocated array of \texttt{n*n} doubles satisfying the invariants (diagonal zero, symmetric, non-negative).
  \item \texttt{-1} if \texttt{malloc} fails. \texttt{inst->dist} is left \texttt{NULL}.
\end{itemize}

\textbf{What it does:}
\begin{enumerate}[label=\arabic*.]
  \item Allocate \texttt{n*n} doubles. Return \texttt{-1} on failure.
  \item For all $i$: set \texttt{DIST(inst,i,i) = 0.0}.
  \item For all $i < j$: compute $d = \sqrt{(x_i - x_j)^2 + (y_i - y_j)^2}$ using \texttt{sqrt()} from \texttt{<math.h>}. Set both \texttt{DIST(inst,i,j) = d} and \texttt{DIST(inst,j,i) = d}.
  \item Return \texttt{0}.
\end{enumerate}

\textbf{Why only the upper triangle?} Computing only $j > i$ means each pair is computed once. Mirroring enforces $D_{ij} = D_{ji}$ \emph{by construction} from the same value, not from two independent computations that might differ by a floating-point epsilon.

\textbf{The \texttt{sqrt} vs \texttt{sqrtf} issue:} \texttt{sqrt()} from \texttt{<math.h>} takes and returns \texttt{double}. \texttt{sqrtf()} takes and returns \texttt{float}. Since both coordinates and distances are \texttt{double} in this design, use \texttt{sqrt()}. If you used \texttt{float} and called \texttt{sqrt()} instead of \texttt{sqrtf()}, the argument would be implicitly promoted to \texttt{double} for the computation and then truncated back to \texttt{float} on assignment---not undefined behaviour, but needless and slightly imprecise. Always match the function suffix to your type.

\end{apispec}

\newpage
